\chapter{Introducci\'on y objetivos}
\label{cap:introduccion}

Este cap\'itulo sit\'ua el trabajo en su \'ambito, delimita el problema que aborda y
enuncia los objetivos que gu\'ian la investigaci\'on. Adem\'as, precisa el alcance del
estudio y describe la organizaci\'on de la memoria.

\section{Contexto y motivaci\'on}
\label{sec:contexto}

La manipulaci\'on rob\'otica requiere generar acciones continuas a partir de
observaciones sensoriales de alta dimensi\'on. El aprendizaje por imitaci\'on aborda esa
dificultad ajustando una pol\'itica sobre demostraciones humanas, sin necesidad de
dise\~nar una funci\'on de recompensa \cite{ross2011dagger}.

Sin embargo, las demostraciones humanas son multimodales. Ante una misma observaci\'on,
un operario puede ejecutar acciones distintas y todas ellas v\'alidas. Los m\'etodos de
clonaci\'on de comportamiento basados en regresi\'on promedian esas alternativas y
producen trayectorias incoherentes \cite{florence2021ibc}.

Para evitar ese promediado, Chi \textit{et al.} \cite{chi2025diffusion} proponen
\textit{Diffusion Policy}, que representa la acci\'on como un proceso de difusi\'on
condicionado por la observaci\'on. El m\'etodo hereda as\'i la capacidad de los modelos
generativos de difusi\'on para representar distribuciones multimodales
\cite{ho2020ddpm}.

En su variante visuomotora, \textit{Diffusion Policy} se compone de dos bloques. Un
codificador visual transforma las im\'agenes en un vector de caracter\'isticas, y una red
generativa produce la secuencia de acciones condicionada por dicho vector. Por tanto, la
representaci\'on visual determina qu\'e informaci\'on alcanza al bloque generativo.

\section{Problema abordado}
\label{sec:problema}

La investigaci\'on sobre \textit{Diffusion Policy} ha modificado principalmente el
bloque generativo y la modalidad de entrada. Algunos trabajos sustituyen las im\'agenes
por nubes de puntos \cite{ze2024dp3}, mientras que otros reducen el coste de inferencia
\cite{prasad2024consistency}.

El trabajo original incluye una ablaci\'on de ResNet y CLIP en la tarea \textit{Square}
\cite{chi2025diffusion}. Adem\'as, un estudio posterior compara ResNet-18 y DINOv3 en
cuatro tareas, incluida Push-T \cite{egbe2025dinov3dp}. Sin embargo, ninguno contrasta
conjuntamente ResNet-18, DINOv2 y CLIP bajo el mismo protocolo de Push-T.

De ah\'i la pregunta que gu\'ia el trabajo: dentro de una misma \textit{Diffusion Policy}
entrenada sobre un conjunto reducido de demostraciones de Push-T, \textquestiondown qu\'e diferencia
de rendimiento y de coste separa a cinco bloques visuales concretos, y con qu\'e grado de
certeza puede afirmarse esa diferencia?

Conviene precisar desde el principio en qu\'e nivel se formula esa pregunta. Cada bloque
visual se lleva a un \'unico artefacto entrenado, es decir, a un punto de control
concreto obtenido con una sola ejecuci\'on de entrenamiento. La comparaci\'on que este
trabajo puede sostener es, por tanto, la comparaci\'on de esos cinco artefactos sobre un
conjunto de condiciones iniciales, y no la estimaci\'on del rendimiento esperado al
repetir cada estrategia de entrenamiento. El \cref{sec:alcance} justifica esa restricci\'on
y el \cref{sec:seleccion-prueba} la traslada a la especificaci\'on estad\'istica.

\section{Objetivos}
\label{sec:objetivos}

El objetivo general consiste en cuantificar y comparar el rendimiento en bucle cerrado y
el coste computacional de cinco bloques visuales integrados como codificador de
observaciones de una misma \textit{Diffusion Policy} en Push-T, cada uno entrenado una
sola vez bajo un protocolo documentado, y en delimitar el alcance inferencial de esa
comparaci\'on.

Este prop\'osito se concreta en los objetivos espec\'ificos siguientes:

\begin{enumerate}[label=\arabic*.,leftmargin=*]
  \item Implementar e integrar los cinco bloques visuales dentro de la misma
        \textit{Diffusion Policy}, manteniendo constantes la red generativa, el conjunto
        de datos y la interfaz de condicionamiento.
  \item Estimar, sobre un bloque de condiciones iniciales disjunto y preregistrado, las
        diferencias de rendimiento entre los cinco puntos de control, acompa\~nadas de su
        incertidumbre y con control de la multiplicidad.
  \item Caracterizar el coste computacional de cada variante en cuatro ejes:
        par\'ametros entrenables, tiempo por \'epoca, latencia de inferencia desglosada y
        pico de memoria gr\'afica.
  \item Acotar la validez de la comparaci\'on: cuantificar el sesgo del conjunto de
        selecci\'on y establecer qu\'e no puede concluirse a partir de una \'unica
        ejecuci\'on de entrenamiento por variante.
\end{enumerate}

\section{Alcance del trabajo}
\label{sec:alcance}

El estudio se circunscribe a la tarea Push-T en simulaci\'on, que consiste en desplazar
una pieza con forma de T hasta una posici\'on objetivo mediante un efector circular
\cite{chi2025diffusion}. Todas las variantes comparten la misma pol\'itica generativa, el
conjunto de datos y el protocolo de evaluaci\'on. La variable manipulada es el bloque
visual completo, es decir, la arquitectura del codificador, sus pesos iniciales, su
estrategia de adaptaci\'on y su cadena de preprocesado, que var\'ian de forma conjunta.
Las desviaciones de presupuesto y tama\~no de lote se documentan de forma expl\'icita.

La premisa de dise\~no que m\'as condiciona el alcance es deliberada y se declara aqu\'i:
\textbf{cada variante se entrena una sola vez}, con la semilla \num{42}. No se trata de
una omisi\'on, sino de una decisi\'on de reparto del presupuesto de c\'omputo. Las cinco
ejecuciones suman \SI{237.1}{\hour} sobre la \'unica GPU disponible, seg\'un el
\cref{sec:coste}; replicar la matriz con tres semillas habr\'ia exigido unas
\SI{700}{\hour}, incompatibles con el calendario del trabajo. En consecuencia, la unidad
sobre la que este estudio infiere es el \emph{artefacto entrenado}, no la estrategia de
entrenamiento: los intervalos y contrastes del \cref{cap:resultados} describen c\'omo se
comportan cinco puntos de control fijos ante nuevas condiciones iniciales, y no c\'omo se
comportar\'ia cada estrategia si volviera a entrenarse. El
\cref{sec:seleccion-prueba} formaliza este estimando y el \cref{sec:limitaciones} recoge
sus consecuencias.

Quedan fuera del alcance cuatro aspectos. En primer lugar, la validaci\'on sobre un robot
f\'isico, puesto que el trabajo opera \'unicamente en simulaci\'on. En segundo lugar, el
an\'alisis de robustez frente a perturbaciones visuales. En tercer lugar, la estimaci\'on
de la variabilidad entre ejecuciones de entrenamiento, que exigir\'ia varias semillas por
variante. En cuarto lugar, las ablaciones factoriales que separar\'ian la inicializaci\'on
de los pesos del preprocesado y de la agregaci\'on espacial. Los dos \'ultimos se
concretan como l\'ineas de continuaci\'on en el \cref{sec:trabajo-futuro}.

\section{Estructura de la memoria}
\label{sec:estructura}

La memoria se organiza en seis cap\'itulos. El \cref{cap:estado-del-arte} revisa los
fundamentos de los modelos de difusi\'on, su aplicaci\'on al control visuomotor y las
familias de codificadores visuales consideradas. El \cref{cap:metodologia} describe el
dise\~no experimental, el conjunto de datos, las variantes evaluadas y las m\'etricas.

A continuaci\'on, el \cref{cap:resultados} presenta las curvas de entrenamiento, la
comparativa de rendimiento y el coste computacional de cada variante. El
\cref{cap:conclusiones} interpreta esos resultados, se\~nala las limitaciones del estudio
y propone l\'ineas de continuaci\'on. Por \'ultimo, el \cref{cap:presupuesto} detalla los
costes del proyecto.
